% Insert after the full characteristic-zero proof of PC30 and before
% the paragraph beginning ``Let chi_{c,k}''.  PC30 itself is retained.
\paragraph{The coefficient map of the same collision-jet calculation.}
The nonzero factorial in the preceding proof also determines the
entire coefficient-specialization map.  Fix an ordered tuple
\(\boldsymbol\rho=(\rho_1,\ldots,\rho_k)\) occurring in (PC30),
and retain \(m_i=m_{\rho_i}\), \(\lambda=\sum_i\rho_i\),
\(D=\sum_i(m_i-1)\), and the original coordinates
\(x_i=s_i-\rho_i\).  On the actual retained coefficient ring
\(A\subset\mathbb C\) carrying these coordinates and the full
Taylor unit and inverse, the local map of (PC29) is
\[
\begin{aligned}
 \mathcal E_A&=A[x_1,\ldots,x_k]/(x_i^{m_i})_{i=1}^k,
 &J&=\sum_i x_i,\qquad S=\lambda+J,\\
 \mathcal C_A&=A[S]/((S-\lambda)^{D+1}),
 &\eta_{\boldsymbol\rho}[P]&=U_{\boldsymbol\rho}P(\lambda+J),
 \qquad U_{\boldsymbol\rho}=\prod_i j_{\rho_i}(\upsilon_h).
\end{aligned}
\tag{PC30a}\label{eq:pc30a-update}
\]
Every original coefficient of \(j_{\rho_i}(\upsilon_h)\),
including its inverse, is retained.  This is a module map for
the displayed original sum coordinate, since multiplication
by \(S\) on the source is multiplication by \(\lambda+J\)
on the target.  Its underlying unweighted ring map has value
one at one, and its weighted map has the original value
\(U_{\boldsymbol\rho}\) at one.

For the specified coefficient map \(\alpha:A\to k_0\) where
it is defined, with finite target of characteristic
\(p>\max_i m_i\), put \(\mathscr R=A_{\ker\alpha}\).
The image of \(\alpha\) is a finite field, so its kernel is
maximal, and the map extends by
\(a/b\mapsto\alpha(a)/\alpha(b)\).  The characteristic-zero
domain \(\mathscr R\) has every integer prime to \(p\) as
a unit and \(p\) as a nonzero nonunit.  Thus all the following
maps are over this specified coefficient localization, with
no extension to a larger CRT stage being presumed.

For every \(0\leq r\leq D\) retain the complete vector
\[
 V_r=\sum_{\substack{0\leq b_i<m_i\\\sum_i b_i=r}}
               \frac{x^{\mathbf b}}{\prod_i b_i!},\quad
 W_r=U_{\boldsymbol\rho}V_r,
 \qquad J^r=r!V_r,
 \qquad JW_r=(r+1)W_{r+1}.
\tag{PC30b}\label{eq:pc30b-update}
\]
Here \(W_{D+1}=0\).  Each denominator is a unit in
\(\mathscr R\).  To verify the first multiplication identity,
expand \((\sum_i x_i)^r\): the surviving monomial
\(x^{\mathbf b}\) has exactly \(r!/\prod_i b_i!\)
ordered assignments.  For the second, its degree-\(r+1\)
coefficient is
\(\sum_{i:b_i>0}b_i/\prod_i b_i!=(r+1)/\prod_i b_i!\).
Multiplication by the original unit proves the formula for
\(W_r\).  A tuple of every degree \(r\leq D\) exists by
successively filling the capacities \(m_i-1\); choose one
pivot \(\mathbf b^*(r)\) for each such degree.  The functionals
\(\lambda_r(x)=\mathbf b^*(r)!\,[x^{\mathbf b^*(r)}]x\)
satisfy \(\lambda_r(V_t)=\delta_{rt}\).  Consequently
\(\pi(x)=\sum_r\lambda_r(x)V_r\) is a projection onto
\(\Lambda=\bigoplus_{r=0}^D\mathscr R V_r\).
Subtracting this projection leaves exactly the span \(F\)
of the nonpivot monomials.  Conjugating by the entire unit
therefore proves
\[
\begin{gathered}
 \mathcal E_{\mathscr R}=\Lambda_U\oplus U_{\boldsymbol\rho}F,
 \qquad \Lambda_U=\bigoplus_{r=0}^D\mathscr R W_r,\\
 L_U:=\eta_{\boldsymbol\rho}(\mathcal C_{\mathscr R})
       =\bigoplus_{r=0}^D r!\mathscr R W_r,
 \qquad
 \eta_{\boldsymbol\rho}((S-\lambda)^r)=r!W_r,\\
 \Lambda_U/L_U\simeq\bigoplus_{r=0}^D\mathscr R/(p^{d_r}),
 \quad d_r=\sum_{a\geq1}\left\lfloor r/p^a\right\rfloor,
 \quad [\sum_r c_rW_r]\longmapsto(c_r\bmod p^{d_r})_r.
\end{gathered}
\tag{PC30c}\label{eq:pc30c-update}
\]
Counting the factors of \(p\) in each integer up to \(r\)
gives \(r!=p^{d_r}\varepsilon_r\) with
\(\varepsilon_r\in\mathscr R^\times\), proving the quotient
formula.  Its kernel and surjectivity follow in the displayed
basis.  In particular the cokernel into the full local tensor
also retains the free summand \(U_{\boldsymbol\rho}F\).
The containing lattice is the exact saturation:
\((L_U\otimes\operatorname{Frac}\mathscr R)\cap
\mathcal E_{\mathscr R}=\Lambda_U\).  Indeed every factorial
becomes invertible in the fraction field, and the free
complement in the displayed splitting has zero component
for an element in that rational span.  No original vector
has been divided by its factorial inside \(\mathcal E_{\mathscr R}\).

Apply the actual map to \(B=k_0\), or the quotient map to
\(B=\mathscr R/p\mathscr R\), and mark base changes by bars.
All pivot maps and the full inverse unit base change, so the
\(\overline W_r\) remain independent.  The exact result is
\[
\begin{gathered}
 \overline\eta_{\boldsymbol\rho}((S-\overline\lambda)^r)
       =\overline{r!}\,\overline W_r,\qquad
 \ker\overline\eta_{\boldsymbol\rho}
       =((S-\overline\lambda)^p),\\
 \operatorname{im}\overline\eta_{\boldsymbol\rho}
       =\bigoplus_{r=0}^{\min(D,p-1)}B\overline W_r,
 \qquad
 \operatorname{ind}(M_{S-\overline\lambda}
          |\overline{\mathcal C}_{\mathscr R})=D+1,
 \qquad
 \operatorname{ind}(M_{\overline J}
          |\operatorname{im}\overline\eta_{\boldsymbol\rho})
       =\min(p,D+1).
\end{gathered}
\tag{PC30d}\label{eq:pc30d-update}
\]
The kernel ideal is zero if \(D<p\).  The factorial is a
unit below \(p\) and zero from \(p\) onwards, proving the
image and kernel.  In characteristic \(p\),
\(\overline J^p=\sum_i\overline x_i^{p}=0\), whereas
\(\overline J^{\min(p-1,D)}\overline U_{\boldsymbol\rho}
=\overline{\min(p-1,D)!}\,
\overline W_{\min(p-1,D)}\ne0\).  These identities prove
the two different indices with their exact module targets.
They are joined by the explicit free resolution and its
connecting map
\[
\begin{gathered}
 0\longrightarrow\mathcal C_{\mathscr R}
  \xrightarrow{\eta_{\boldsymbol\rho}}\Lambda_U
       \longrightarrow\Lambda_U/L_U\longrightarrow0,\\
 \operatorname{Tor}_1^{\mathscr R}(\Lambda_U/L_U,B)
       =\bigoplus_{r=p}^D B\epsilon_r,
 \qquad \partial\epsilon_r=(S-\overline\lambda)^r
       \in\ker\overline\eta_{\boldsymbol\rho}.
\end{gathered}
\tag{PC30e}\label{eq:pc30e-update}
\]
In the retained bases its differential is
\(\operatorname{diag}(0!,\ldots,D!)\).  Its base change
has exactly the displayed kernel, proving this Tor group
and its connecting map, with differential sign positive.

There is also an exact formula for collisions after a
coefficient map actually defined on the full one-factor
CRT data in (PC28)--(PC30), with \(p>m_\rho\) for every
\(\rho\in Z\).  Retain those idempotents and
their original inverses, so their base changes still give
the product of the ordered local factors.  Let
\(\mathcal Z_{k,p}=\{\alpha(\sum_i\rho_i):
                         \boldsymbol\rho\in Z^k\}\)
be the set of attained residual sums.  Define
\[
 \psi_{h,k}(X)=\prod_{\mu\in\mathcal Z_{k,p}}
       (X-\mu)^{n_\mu},\qquad
 n_\mu=\max_{\substack{\boldsymbol\rho\in Z^k\\
                         \alpha(\sum_i\rho_i)=\mu}}
           \min\left(p,1+\sum_i(m_{\rho_i}-1)\right).
\tag{PC30f}\label{eq:pc30f-update}
\]
Every maximum in this display is over a nonempty set.
When \(Z\) is empty, \(k\geq1\) makes the tuple set
empty, so the product is \(\psi_{h,k}=1\); both the
original cyclic source \(k_0[X]/(1)\) and its tensor
image are zero.  The global inequality on the orders
is vacuous in this empty case.
The multiplication matrix on each factor is
\(\mu+\overline J\), whose least annihilating power
is (PC30d).  More explicitly, write a nonzero polynomial
as \(P(X)=(X-\mu)^eQ(X)\) with \(Q(\mu)\ne0\).
The matrix \(Q(\mu+\overline J)\) is the nonzero scalar
\(Q(\mu)\) times \(I+H\), with \(H\) a polynomial
in \(\overline J\) having zero constant term and hence
nilpotent.  Its inverse is the finite geometric series
\(Q(\mu)^{-1}\sum_{j\geq0}(-H)^j\), stopped after
the nilpotence index.  Thus \(P(\mu+\overline J)=0\)
exactly when \(\overline J^e=0\).  This proves that
the least monic annihilating polynomial is the indicated
power of \(X-\mu\).  On a product, a polynomial annihilates the
operator exactly when it annihilates every factor;
its monic annihilator is their least common multiple.
Taking the maximum exponent at each equal residual sum
proves (PC30f).  The original reduced polynomial
\(\overline\chi_{h,k}\) still annihilates all factors,
so \(\psi_{h,k}\mid\overline\chi_{h,k}\).  The original
identity vector generates the polynomial image, and
the full reduced unit is invertible; hence the precise
specialized map of (PC29) has
\[
 \ker\overline\eta_k
       =(\psi_{h,k})/(\overline\chi_{h,k}),\qquad
 \operatorname{im}\overline\eta_k
       \simeq k_0[X]/(\psi_{h,k}),
 \quad [P]\longmapsto\overline U_kP(\overline S).
\tag{PC30g}\label{eq:pc30g-update}
\]
The type of this isomorphism is a module isomorphism;
the unweighted map has the corresponding algebra
isomorphism.  Its kernel is the annihilator ideal
just proved, and its image is the definition of that
polynomial image, proving both statements.  This
retains new residual collisions instead of reducing
the characteristic-zero maximum without its map.
It makes no claim that a local coefficient map extends
to all the sum-CRT coefficients: the actual separating
idempotent at degree \(k\geq p\) in the original
quartet corner gives \(p^{-1}\) by its divided
difference (BTF.9), and forbids such an extension.
The characteristic-zero formula (PC30), its critical
tuple quotient below, and the completed complex
comparison retain their original domains throughout.
